\chapter{The winding-number argument}
\label{chap:winding}

% archon:covers Gluck/Winding.lean

This is the analytic core of Gluck's proof. We vary the four breakpoints of the
bicircle (\cref{chap:bicircle}) over a concrete convex $2$-disk $D$ of
configurations, show that the error vector $E$ traces a loop winding once around
the origin as the configuration runs over $\partial D$, and invoke a planar
degree theorem: a continuous map on a disk whose boundary winds once around $0$
has an interior zero. That interior zero is a configuration whose bicircle closes
up.

\medskip

\noindent\emph{Mathlib gap.} Mathlib provides holomorphic winding (Cauchy
integral formula) but \emph{no} topological planar winding number of a continuous
loop, and no ``nonzero boundary winding $\Rightarrow$ interior zero'' theorem.
Both are built project-side here; see \cref{thm:existence_of_zero}. We use a
finite-dimensional $2$-disk of breakpoint configurations, \emph{not} loops in the
diffeomorphism group $\mathrm{Diff}(S^1)$ (which has no Mathlib homotopy
infrastructure): this is the strategic replacement for the source's
``contractible loop of diffeomorphisms''.

\section{Topological winding-number primitives}

Mathlib has no topological winding number of a continuous loop $S^1 \to
\mathbb{C}\setminus\{0\}$. We build one project-side from the exponential covering
$\mathrm{Circle.exp} : \mathbb{R} \to S^1$ (Mathlib's
\texttt{Circle.isCoveringMap\_exp}) by lifting and measuring the total angle
increment. These are the supporting declarations that
\cref{thm:existence_of_zero} and \cref{lem:error_map_winds_boundary} consume; each
carries a \texttt{\textbackslash lean} pin so the Lean file is in $1$-to-$1$ correspondence with the
blueprint.

\begin{definition}

\leanok
[Angle lift of a circle loop]
  \label{def:angle_lift}
  % Lean (private, not pinned): angleLift
  Project-bespoke. The continuous \emph{angle lift} of a loop $g : [0,1] \to S^1$:
  the path $\widetilde g : [0,1] \to \mathbb{R}$ obtained by lifting $g$ along the
  covering $\mathrm{Circle.exp}$ (Mathlib \texttt{IsCoveringMap.liftPath}),
  starting at a chosen real preimage of $g(0)$.
\end{definition}

\begin{lemma}

\leanok
[The angle lift lifts]
  \label{lem:angle_lift_lifts}
  % Lean (private, not pinned): angleLift_lifts
  \uses{def:angle_lift}
  Project-bespoke. $\mathrm{Circle.exp}(\widetilde g(t)) = g(t)$ for all $t$.
\end{lemma}
\begin{proof}

\leanok
  Immediate from the lifting property \texttt{IsCoveringMap.liftPath\_lifts}.
\end{proof}

\begin{definition}

\leanok
[Winding number of a circle loop]
  \label{def:winding_number}
  % Lean (private, not pinned): windingNumber
  \uses{def:angle_lift}
  Project-bespoke. The \emph{winding number} of $g : [0,1] \to S^1$ is the
  normalised total angle increment $W(g) = (\widetilde g(1) - \widetilde g(0))/(2\pi)$.
\end{definition}

\begin{lemma}

\leanok
[A constant loop does not wind]
  \label{lem:winding_number_const}
  % Lean (private, not pinned): windingNumber_const
  \uses{def:winding_number}
  Project-bespoke. $W$ of a constant loop is $0$.
\end{lemma}
\begin{proof}

\leanok
  The lift of a constant loop is constant (\texttt{liftPath\_const}), so the
  increment is $0$.
\end{proof}

\begin{lemma}

\leanok
[Integer-valued continuous functions on $[0,1]$ are constant]
  \label{lem:int_valued_eq}
  % Lean (private, not pinned): int_valued_eq
  Project-bespoke. A continuous $q : [0,1] \to \mathbb{R}$ taking only integer
  values is constant.
\end{lemma}
\begin{proof}

\leanok
  If $q(a) \ne q(b)$, the intermediate value theorem (on the connected $[0,1]$)
  forces $q$ to take the non-integer value halfway between two consecutive
  integers, a contradiction.
\end{proof}

\begin{lemma}

\leanok
[Any lift computes the winding number]
  \label{lem:winding_number_eq_div_of_lift}
  % Lean (private, not pinned): windingNumber_eq_div_of_lift
  \uses{def:winding_number, lem:int_valued_eq, lem:angle_lift_lifts}
  Project-bespoke. If $\varphi : [0,1] \to \mathbb{R}$ is \emph{any} continuous
  lift of $g$ (i.e.\ $\mathrm{Circle.exp}\circ\varphi = g$), then $W(g) =
  (\varphi(1) - \varphi(0))/(2\pi)$.
\end{lemma}
\begin{proof}
  \uses{lem:int_valued_eq}
  \leanok
  $\varphi - \widetilde g$ is continuous and $2\pi\mathbb{Z}$-valued (two lifts of
  the same loop differ by a constant in the fibre, via \texttt{Circle.exp\_eq\_exp}),
  hence constant by \cref{lem:int_valued_eq}; so the two increments agree.
\end{proof}

\begin{lemma}

\leanok
[Winding number is a free-homotopy invariant]
  \label{lem:winding_number_eq_of_homotopy}
  % Lean (private, not pinned): windingNumber_eq_of_homotopy
  \uses{def:winding_number, lem:winding_number_eq_div_of_lift, lem:int_valued_eq}
  Project-bespoke. If $H : [0,1]\times[0,1] \to S^1$ is a continuous homotopy
  \emph{through loops} (each slice $H(s,\cdot)$ is a loop), then $W(H(0,\cdot)) =
  W(H(1,\cdot))$.
\end{lemma}
\begin{proof}
  \uses{lem:winding_number_eq_div_of_lift, lem:int_valued_eq}
  \leanok
  Lift $H$ along $\mathrm{Circle.exp}$ to $\widetilde H$ (Mathlib
  \texttt{IsCoveringMap.liftHomotopy}). For each $s$, $(\widetilde H(s,1) -
  \widetilde H(s,0))/2\pi$ equals $W(H(s,\cdot))$ (by
  \cref{lem:winding_number_eq_div_of_lift}) and is integer-valued (loop condition
  $H(s,0)=H(s,1)$), hence constant in $s$ by \cref{lem:int_valued_eq}.
\end{proof}

\begin{lemma}

\leanok
[Winding number is additive under pointwise product]
  \label{lem:winding_number_mul}
  % Lean (private, not pinned): windingNumber_mul
  \uses{def:winding_number, lem:winding_number_eq_div_of_lift, lem:angle_lift_lifts}
  Project-bespoke. For circle loops $g, h$, $W(g\cdot h) = W(g) + W(h)$.
\end{lemma}
\begin{proof}

\leanok
  $\widetilde g + \widetilde h$ is a continuous lift of $g\cdot h$ (since
  $\mathrm{Circle.exp}(\widetilde g + \widetilde h) = \mathrm{Circle.exp}\,
  \widetilde g \cdot \mathrm{Circle.exp}\,\widetilde h = g\cdot h$ by
  \texttt{Circle.exp\_add} and \cref{lem:angle_lift_lifts}); apply
  \cref{lem:winding_number_eq_div_of_lift} and split the increment.
\end{proof}

\begin{definition}

\leanok
[Radial projection to the circle]
  \label{def:circle_proj}
  % Lean (private, not pinned): circleProj
  Project-bespoke. For $z \ne 0$, $\mathrm{circleProj}(z) = z/\lVert z\rVert \in
  S^1$, the radial normalisation onto the unit circle.
\end{definition}

\begin{lemma}

\leanok
[Radial projection formula]
  \label{lem:circle_proj_eq}
  % Lean (private, not pinned): circleProj_eq
  \uses{def:circle_proj}
  Project-bespoke. As a complex number, $\mathrm{circleProj}(z) = z/\lVert z\rVert$.
\end{lemma}
\begin{proof}

\leanok
  Unfolding the definition; membership in $S^1$ is $\lVert z/\lVert z\rVert\rVert = 1$.
\end{proof}

\begin{lemma}

[Radial projection: congruence, coercion, multiplicativity]
  \label{lem:circle_proj_helpers}
  % Lean (private, not pinned): circleProj_mul, circleProj_coe, circleProj_congr
  \uses{def:circle_proj, lem:circle_proj_eq}
  Project-bespoke (private in-file helpers). (i) \emph{congruence}
  ($\mathrm{circleProj\_congr}$): $a=b \Rightarrow \mathrm{circleProj}(a) =
  \mathrm{circleProj}(b)$; (ii) \emph{coercion} ($\mathrm{circleProj\_coe}$): for
  $z\in S^1$, $\mathrm{circleProj}(z) = z$; (iii) \emph{multiplicativity}
  ($\mathrm{circleProj\_mul}$): $\mathrm{circleProj}(c\,z) =
  \mathrm{circleProj}(c)\,\mathrm{circleProj}(z)$ for $c,z\ne 0$.
\end{lemma}
\begin{proof}

  All three by unfolding \cref{lem:circle_proj_eq} ($\mathrm{circleProj}(z) =
  z/\lVert z\rVert$): (i) is substitution; (ii) uses $\lVert z\rVert = 1$ on
  $S^1$; (iii) uses $\lVert c z\rVert = \lVert c\rVert\,\lVert z\rVert$.
\end{proof}

\begin{definition}

\leanok
[Normalised loop of a $\mathbb{C}$-loop]
  \label{def:norm_loop}
  % Lean (private, not pinned): normLoop
  \uses{def:circle_proj}
  Project-bespoke. For a nonvanishing loop $\gamma : [0,1] \to \mathbb{C}$,
  $\mathrm{normLoop}\,\gamma$ is the circle loop $t \mapsto
  \mathrm{circleProj}(\gamma(t))$.
\end{definition}

\begin{definition}

\leanok
[Winding number of a $\mathbb{C}$-loop]
  \label{def:winding_number_c}
  \lean{Gluck.windingNumberC}
  \uses{def:winding_number, def:norm_loop}
  Project-bespoke. The winding number of a nonvanishing loop $\gamma : [0,1] \to
  \mathbb{C}$ is $W_{\mathbb{C}}(\gamma) = W(\mathrm{normLoop}\,\gamma)$.
\end{definition}

\begin{lemma}

\leanok
[Scaling invariance of the $\mathbb{C}$-winding number]
  \label{lem:winding_number_c_const_mul}
  % Lean (private, not pinned): windingNumberC_const_mul
  \uses{def:winding_number_c, def:norm_loop, lem:winding_number_mul,
        lem:winding_number_const, lem:circle_proj_helpers}
  Project-bespoke. For $c\ne 0$ and a nonvanishing loop $\gamma$,
  $W_{\mathbb{C}}(c\cdot\gamma) = W_{\mathbb{C}}(\gamma)$.
\end{lemma}
\begin{proof}

\leanok
  By \cref{lem:circle_proj_helpers}(iii), $\mathrm{normLoop}(c\cdot\gamma) =
  \mathrm{circleProj}(c)\cdot\mathrm{normLoop}\,\gamma$ as a product of circle
  loops (the first factor constant); apply \cref{lem:winding_number_mul} and
  \cref{lem:winding_number_const} ($W$ of the constant factor is $0$).
\end{proof}

\begin{lemma}

\leanok
[Positive-scalar-field invariance of the $\mathbb{C}$-winding number]
  \label{lem:winding_number_c_pos_scalar_field}
  \lean{Gluck.windingNumberC_posScalarField}
  \uses{def:winding_number_c, def:norm_loop, lem:winding_number_eq_of_homotopy}
  Project-bespoke. Let $\gamma:[0,1]\to\mathbb{C}$ be a nowhere-zero loop and
  $c:[0,1]\to\mathbb{R}$ continuous with $c(t)>0$ for all $t$ (and $c$ a loop,
  $c(0)=c(1)$). Then the scaled loop $t\mapsto c(t)\,\gamma(t)$ is nowhere zero and
  $W_{\mathbb{C}}(c\cdot\gamma)=W_{\mathbb{C}}(\gamma)$. This generalises
  \cref{lem:winding_number_c_const_mul} from a constant scalar to a positive scalar
  \emph{field}; it is what lets a positive configuration-dependent prefactor
  $c(z)=1/\lambda(z)$ be stripped from the clean arc-length error map
  (\cref{lem:clean_error_winds}).
\end{lemma}
\begin{proof}
  \uses{def:winding_number_c, lem:winding_number_eq_of_homotopy}
  \leanok
  The straight-line scalar homotopy
  $H(u,t)=\bigl((1-u)+u\,c(t)\bigr)\,\gamma(t)$ connects $\gamma$ ($u=0$) to
  $c\cdot\gamma$ ($u=1$). Its scalar factor $(1-u)+u\,c(t)$ is a convex combination
  of $1>0$ and $c(t)>0$, hence $>0$ for all $(u,t)\in[0,1]^2$; since
  $\gamma(t)\neq0$, $H$ is nowhere zero. $H$ is continuous and, being a real-scalar
  multiple of the loop $\gamma$, a loop in $t$ at each $u$. Normalising and applying
  \cref{lem:winding_number_eq_of_homotopy} gives
  $W_{\mathbb{C}}(c\cdot\gamma)=W_{\mathbb{C}}(\gamma)$. Nowhere-vanishing of
  $c\cdot\gamma$ is the $u=1$ slice.
\end{proof}

\begin{lemma}

\leanok
[Winding is stable under small perturbations]
  \label{lem:winding_number_c_perturb}
  \lean{Gluck.windingNumberC_eq_of_perturb}
  \uses{def:winding_number_c, def:norm_loop, lem:winding_number_eq_of_homotopy}
  Project-bespoke. Let $\gamma,\gamma'$ be nowhere-zero loops with
  $\lVert\gamma'(t)-\gamma(t)\rVert < \lVert\gamma(t)\rVert$ for all $t$. Then
  $W_{\mathbb{C}}(\gamma') = W_{\mathbb{C}}(\gamma)$.
\end{lemma}
\begin{proof}

\leanok
  The straight-line homotopy $\gamma_\tau = \gamma + \tau(\gamma'-\gamma)$ stays
  nowhere zero on $[0,1]^2$ by the reverse triangle inequality
  ($\lVert\gamma_\tau\rVert \ge \lVert\gamma\rVert - \lVert\gamma'-\gamma\rVert >
  0$); normalise it and apply \cref{lem:winding_number_eq_of_homotopy}.
\end{proof}

\begin{definition}

\leanok
[Boundary loop of the unit disk]
  \label{def:disk_boundary_loop}
  \lean{Gluck.diskBoundaryLoop}
  Project-bespoke. For $F : \overline{\mathbb{D}} \to \mathbb{C}$, the
  \emph{boundary loop} is $t \mapsto F(e^{2\pi i t})$.
\end{definition}

\begin{lemma}

\leanok
[Boundary loop is nonvanishing]
  \label{lem:disk_boundary_loop_ne_zero}
  \lean{Gluck.diskBoundaryLoop_ne_zero}
  \uses{def:disk_boundary_loop}
  Project-bespoke. If $F \ne 0$ on $\partial\mathbb{D}$, the boundary loop is
  nowhere zero.
\end{lemma}
\begin{proof}

\leanok
  Each point $e^{2\pi i t}$ lies on $\partial\mathbb{D}$, where $F \ne 0$.
\end{proof}

% SOURCE: [DeTurck-Gluck], §6, p. 12-13, lines 218-237 (read from references/deturck-gluck-fourvertex.txt)
% SOURCE QUOTE: "The picture has eight portions, located at the eight points of
% the compass. Each contains a small circle which shows a curvature step function
% κ0 ◦ h , and a larger curve which realizes it ... The figures are arranged so
% that the one in the east has an error vector pointing east, the one in the
% northeast has an error vector pointing northeast, and so on. The eight
% diffeomorphisms h are samples from a loop of diffeomorphisms of the circle, and
% the corresponding error vectors E(h) wind once around the origin. The loop is
% contractible in Diff(S1) because each diffeomorphism leaves the bottom point of
% the circle fixed."
\begin{definition}

\leanok
[Configuration disk]
  \label{def:configuration_space}
  \lean{Gluck.configSpace}
  \uses{def:four_arc_step_function}
  \textit{Source: DeTurck--Gluck, §6 (the loop of diffeomorphisms, finitised).}
  Fix $0 < a < b$ and a radius $0 < \delta \le \pi/8$. The \emph{configuration
  disk} is the explicit two-parameter family parametrised by the closed unit disk
  $\overline{\mathbb{D}} \subset \mathbb{R}^2$ (coordinates $(x,y)$, $x^2+y^2 \le
  1$): the two leading breakpoints vary while the trailing two are pinned,
  \[
    \mathrm{configSpace}(x,y) \;=\;
      \bigl(\theta_1,\theta_2,\theta_3,\theta_4\bigr)
      \;=\;\Bigl(\tfrac{\pi}{4} + \delta x,\;
                 \tfrac{3\pi}{4} + \delta y,\;
                 \tfrac{5\pi}{4},\;\tfrac{7\pi}{4}\Bigr).
  \]
  The centre $(x,y)=(0,0)$ is the canonical equally-spaced bicircle (breakpoints
  $\tfrac{\pi}{4},\tfrac{3\pi}{4},\tfrac{5\pi}{4},\tfrac{7\pi}{4}$, four arcs of
  length $\pi/2$), which closes up: its chord sum is $V = (e^{3\pi i/4} - e^{\pi
  i/4}) + (e^{7\pi i/4} - e^{5\pi i/4}) = -\sqrt2 + \sqrt2 = 0$; since the bicircle
  error vector is a nonzero scalar multiple of this chord sum
  (\cref{lem:error_vector_of_bicircle}), the canonical bicircle has zero error
  vector. The
  bound $\delta \le \pi/8$ keeps the order constraint $0 < \theta_1 < \theta_2 <
  \theta_3 < \theta_4 < \theta_1 + 2\pi$ on all of $\overline{\mathbb{D}}$ (since
  $|\delta x|, |\delta y| \le \delta \le \pi/8 < \pi/4$). $\mathrm{configSpace}$ is
  affine, hence continuous, and is a homeomorphism onto its (convex) image.
\end{definition}

\begin{lemma}


\leanok
[Configuration disk: order]
  \label{lem:config_space_props}
  % Lean (private, not pinned): configSpace_ordered
  \uses{def:configuration_space}
  Project-bespoke. For $0<\delta\le\pi/8$ and $\lvert x\rvert,\lvert y\rvert\le 1$,
  the order constraint $0<\theta_1<\theta_2<\theta_3<\theta_4<\theta_1+2\pi$ holds
  ($\mathrm{configSpace\_ordered}$).
\end{lemma}
\begin{proof}


\leanok
  Substitute the formulas $\theta_1=\pi/4+\delta x$, etc.; $\lvert\delta
  x\rvert,\lvert\delta y\rvert\le\delta\le\pi/8<\pi/4$ keeps the four breakpoints
  strictly ordered within one period (\texttt{nlinarith}/\texttt{linarith}).
\end{proof}

\begin{remark}
  The source realises the family as samples from a contractible loop of
  circle diffeomorphisms fixing the bottom point; the eight compass pictures are
  eight samples whose error vectors point in the eight compass directions. The
  finite-dimensional disk $D$ records exactly the data that the error vector
  depends on (the four breakpoints), discarding the rest of the diffeomorphism.
  Contractibility of the source's loop corresponds to $D$ being a disk (its
  boundary $\partial D$ is the loop, contractible in $D$).
\end{remark}

\begin{definition}


\leanok
[The error map and its closed form]
  \label{def:error_map}
  \lean{Gluck.errorMap}
  % Lean (private, not pinned): errorMap_eq, continuousOn_errorMap
  \uses{def:configuration_space, def:bicircle_error_vector}
  Project-bespoke. $\mathrm{errorMap}(a,b,\delta)(z) := E_{\mathrm{bi}}\bigl(a,b,
  \mathrm{configSpace}(\delta)(z.\mathrm{re},z.\mathrm{im})\bigr)$ is the bicircle
  error vector of the configuration parametrised by $z\in\overline{\mathbb{D}}$.
  $\mathrm{errorMap\_eq}$ records its closed form $E = s\cdot(V+\sqrt2)$ on the
  disk (from \cref{lem:error_vector_of_bicircle}), and
  $\mathrm{continuousOn\_errorMap}$ its continuity there.
\end{definition}

\begin{definition}


\leanok
[Reverse once-around loop primitives]
  \label{def:reverse_loop}
  \lean{Gluck.windingNumberC_congr}
  % Lean (private, not pinned): negStandardLoop, windingNumber_negStandard, negCircleExpLoop, negCircleExpLoop_ne, windingNumberC_negCircleExp
  \uses{def:winding_number, def:winding_number_c, def:norm_loop,
        lem:winding_number_eq_div_of_lift, lem:circle_proj_helpers}
  Project-bespoke reverse-orientation once-around loops:
  $\mathrm{negStandardLoop}(t)=\mathrm{Circle.exp}(-2\pi t)$, the circle loop
  $t\mapsto e^{-2\pi i t}$, with winding $-1$ (computed directly from the lift
  $\varphi(t)=-2\pi t$ via \cref{lem:winding_number_eq_div_of_lift}); its
  $\mathbb{C}$-valued form
  $\mathrm{negCircleExpLoop}$ (nonvanishing) with $W_{\mathbb{C}}=-1$, and
  $\mathrm{windingNumberC\_congr}$ (pointwise-equal $\mathbb{C}$-loops have equal
  winding). These realise the linear model $L=\delta e^{-i\pi/4}e^{-2\pi i t}$ in
  the proof of \cref{lem:error_map_winds_boundary}.
\end{definition}

% SOURCE: [DeTurck-Gluck], §4 and §6, lines 181-188, 227-237 (read from references/deturck-gluck-fourvertex.txt)
% SOURCE QUOTE: "the corresponding error vectors E(h) will wind once around the
% origin. Furthermore, we will do this so that the loop is contractible in the
% group Diff(S 1) of diffeomorphisms of the circle, and conclude that for some h
% \"inside\" this loop, the curve αh will have error vector E(h) = 0 ."
\begin{lemma}


\leanok
[The error map winds nontrivially on the boundary]
  \label{lem:error_map_winds_boundary}
  \lean{Gluck.errorMap_winding_eq_one}
  % Lean (private, not pinned): errorMap_order, Vpart, Lpart, Lpart_norm, pert_lt, Vpart_ne, remainder_identity, remainder_norm_le, continuous_Vpart_boundary, errVloop, errSVloop, errLloop, errLloop_eq_Lpart
  \uses{def:configuration_space, def:error_map, def:reverse_loop,
        lem:error_vector_of_bicircle, def:winding_number_c}
  \textit{Source: DeTurck--Gluck, §6 (error vectors wind once around the origin).}
  Let $E : \overline{\mathbb{D}} \to \mathbb{C}$ send $(x,y)$ to the bicircle error
  vector of $\mathrm{configSpace}(x,y)$ (\cref{def:configuration_space},
  \cref{lem:error_vector_of_bicircle}). For all sufficiently small $\delta > 0$,
  $E$ is continuous, $E \ne 0$ on the boundary circle $\partial\mathbb{D}$, and the
  boundary loop $E|_{\partial\mathbb{D}}$ has \emph{nonzero} winding number
  $W_{\mathbb{C}}(E|_{\partial\mathbb{D}}) = \pm 1 \ne 0$ about the origin.

  \emph{(Only nonvanishing of the winding number is used downstream
  (\cref{thm:existence_of_zero}); the sign $-1$ here merely reflects the
  orientation of the chosen parametrisation.)}
\end{lemma}
\begin{proof}
  \uses{lem:error_vector_of_bicircle, def:configuration_space, def:winding_number_c,
        lem:winding_number_eq_of_homotopy, def:reverse_loop}
        \leanok
  By \cref{lem:error_vector_of_bicircle}, $E(x,y) = s\cdot V(x,y)$ with the nonzero
  constant scalar $s = \tfrac{1}{ib} - \tfrac{1}{ia}$ (nonzero since $a \ne b$) and
  chord sum
  \[
    V(x,y) = \bigl(e^{i\theta_2} - e^{i\theta_1}\bigr) + \sqrt2,
    \qquad \theta_1 = \tfrac{\pi}{4}+\delta x,\ \theta_2 = \tfrac{3\pi}{4}+\delta y,
  \]
  (using $e^{i\theta_4}-e^{i\theta_3} = e^{7\pi i/4}-e^{5\pi i/4} = \sqrt2$ for the
  pinned trailing breakpoints). $E$ is a composition of continuous maps, hence
  continuous. Multiplying a loop by the fixed nonzero scalar $s$ does not change its
  winding number, so $W_{\mathbb{C}}(E|_{\partial\mathbb{D}}) =
  W_{\mathbb{C}}(V|_{\partial\mathbb{D}})$; it remains to compute the winding of $V$.

  \emph{Linear model.} The first-order Taylor expansion of $V$ at the centre is the
  invertible $\mathbb{R}$-linear map
  \[
    L(x,y) = \tfrac{\delta}{\sqrt2}\bigl[(x-y) - i\,(x+y)\bigr],
  \]
  obtained from $\partial_x V = \delta(-i\,e^{i\pi/4}) = \tfrac{\delta}{\sqrt2}(1-i)$
  and $\partial_y V = \delta(i\,e^{3\pi i/4}) = \tfrac{\delta}{\sqrt2}(-1-i)$. As a
  matrix $L = \delta\,U$ with $U = \tfrac{1}{\sqrt2}\bigl(\begin{smallmatrix} 1 & -1
  \\ -1 & -1\end{smallmatrix}\bigr)$ orthogonal ($\det U = -1$); hence $\lVert
  L(x,y)\rVert =
  \delta$ for every $(x,y)$ on the unit circle, and $V(0,0)=0$ so $V = L + R$ with
  remainder $R$ vanishing to second order: $\lVert R(x,y)\rVert \le C\delta^2$ on
  $\partial\mathbb{D}$ (the map $(x,y) \mapsto e^{i\theta_2}-e^{i\theta_1}$ is smooth
  with second derivatives bounded by $\delta^2$, so the linear-Taylor error is
  $O(\delta^2)$ with an absolute constant $C$).

  \emph{Homotopy to the linear model.} Choose $\delta$ small enough that $C\delta^2
  < \delta$, i.e.\ $\delta < 1/C$ (and $\delta \le \pi/8$ as in
  \cref{def:configuration_space}). Then on $\partial\mathbb{D}$, $\lVert R\rVert \le
  C\delta^2 < \delta = \lVert L\rVert$, so the straight-line homotopy $V_\tau = L +
  \tau R$ ($\tau\in[0,1]$) is nowhere zero on $\partial\mathbb{D}$ (since $\lVert
  V_\tau\rVert \ge \lVert L\rVert - \lVert R\rVert > 0$). By
  \cref{lem:winding_number_eq_of_homotopy} (applied to the induced homotopy of
  normalised circle loops), $W_{\mathbb{C}}(V|_{\partial\mathbb{D}}) =
  W_{\mathbb{C}}(L|_{\partial\mathbb{D}})$.

  \emph{Winding of the linear model.} $L = \delta U$ is an invertible
  $\mathbb{R}$-linear map; the further straight-line homotopy of $U$ to the standard
  reflection through invertible matrices (any invertible matrix is path-connected
  within $\mathrm{GL}_2$ to $\pm$ a rotation, the path stays invertible so the loop
  stays nonzero) carries $L|_{\partial\mathbb{D}}$ to a once-around loop, with the
  reflection $\det U = -1$ fixing the reverse orientation; by \cref{def:reverse_loop}
  (the reverse once-around loop $t\mapsto e^{-2\pi i t}$ has winding $-1$) its winding
  is $-1$. In particular it is nonzero. This is
  the planar analogue of ``a degree-one polynomial winds once''
  (\cref{thm:existence_of_zero}'s FTA analogy).
\end{proof}

% SOURCE: [DeTurck-Gluck], §7, p. 13-14, lines 251-259 (read from references/deturck-gluck-fourvertex.txt)
% SOURCE QUOTE: "Proving something that is obvious by using a robust argument
% reminds us of the proof that the polynomial p0(z) = an zn has a zero by noting
% that it takes any circle in the complex plane to a curve which winds n times
% around the origin. After all, it is obvious that p0(z) has a zero at the
% origin. But this argument is also robust, and applies equally well to the
% polynomial p(z) = an zn + ... + a0 , when z moves around a suitably large
% circle, giving us the usual topological proof of the Fundamental Theorem of
% Algebra."
\begin{theorem}

\leanok
[Nonzero boundary winding forces an interior zero]
  \label{thm:existence_of_zero}
  \lean{Gluck.exists_zero_of_boundary_winding}
  \uses{lem:error_map_winds_boundary}
  \textit{Source: DeTurck--Gluck, §7 (winding-number/FTA principle).}
  Let $F : \overline{\mathbb{D}} \to \mathbb{C}$ be continuous on the closed unit
  $2$-disk, with $F \ne 0$ on the boundary circle $\partial\mathbb{D}$ and the
  boundary loop $F|_{\partial\mathbb{D}}$ having winding number $\ne 0$ about the
  origin. Then $F$ has a zero in the interior of $\mathbb{D}$.
\end{theorem}
\begin{proof}
  \uses{lem:error_map_winds_boundary}
  \leanok
  Suppose, for contradiction, that $F(z) \ne 0$ for all $z \in
  \overline{\mathbb{D}}$. Then $G = F/\lvert F\rvert : \overline{\mathbb{D}} \to
  S^1$ is continuous and extends $g = F|_{\partial\mathbb{D}}/\lvert
  F|_{\partial\mathbb{D}}\rvert$ to the whole disk. The boundary circle
  $\partial\mathbb{D}$ is contractible inside $\overline{\mathbb{D}}$ (shrink to
  the centre), so $g$ is null-homotopic as a map into $S^1$; therefore its degree
  --- which is the winding number of $F|_{\partial\mathbb{D}}$ about $0$ --- is
  zero, contradicting the hypothesis. Hence $F$ vanishes somewhere in
  $\overline{\mathbb{D}}$; the boundary hypothesis $F \ne 0$ on
  $\partial\mathbb{D}$ forces the zero into the interior.

  This is the planar degree principle the source invokes by analogy with the
  topological proof of the Fundamental Theorem of Algebra (\S7): a loop with
  nonzero winding cannot bound a non-vanishing map. Mathlib has no off-the-shelf
  form, so the winding number of a continuous loop $S^1 \to \mathbb{C}\setminus
  \{0\}$, its homotopy invariance, and this null-homotopy argument are built
  project-side (see \cref{chap:winding}'s Lean file).
\end{proof}
